\begin{abstract}
L'integrazione dei \textit{Large Language Model (LLM)} può offrire nuove opportunità per la semplificazione nella gestione e nella manutenzione di infrastrutture server IT.

In particolare, tramite l'emergente protocollo \textit{MCP Server (Model Context Protocol)}, è possibile fornire ad un \textit{LLM} un layer di comunicazione strandardizzato con un sistema IT.
Questo approccio abilita interazioni basate sul linguaggio umano con tutta una serie di vantaggi in termini di semplicità e immediatezza operativa per task che tradizionalmente richiederebbero approcci più impegnativi e complessi.

Tuttavia, è importante considerare la natura aleatoria e non deterministica intrinseca di un sistema basato su AI, pertanto questo lavoro di tesi illustra lo sviluppo di un \textit{server MCP} che permette il relazionarsi di un LLM con un \textit{Hypervisor Proxmox}, con annessa valutazione dei limiti e benefici che tale architettura comporta, ponendo una rigorosa attenzione alla mitigazione dei rischi e alla sicurezza. Attraverso la validazione di questo Proof of Concept, il lavoro dimostra come l'impiego del protocollo \textit{MCP}, unito ad adeguati vincoli architetturali, permetta di governare l'esecuzione di task potenzialmente distruttive in modo controllato e affidabile.
\end{abstract}